% Generated draft: requires mathematical and source review.
\documentclass[11pt]{article}
\usepackage[margin=25mm]{geometry}
\usepackage{fontspec}
\setmainfont{Latin Modern Roman}
\usepackage{amsmath,amssymb,mathtools,unicode-math}
\setmathfont{Latin Modern Math}
\usepackage{enumitem,needspace,xurl,hyperref}
\setlength{\parindent}{0pt}
\setlength{\parskip}{5pt}
\setlength{\emergencystretch}{3em}
\newcommand{\N}{\mathbb{N}}
\newcommand{\Z}{\mathbb{Z}}
\newcommand{\Q}{\mathbb{Q}}
\newcommand{\R}{\mathbb{R}}
\newcommand{\C}{\mathbb{C}}
\newcommand{\F}{\mathbb{F}}
\newcommand{\A}{\mathbb{A}}
\newcommand{\PP}{\mathbb{P}}
\newcommand{\E}{\mathbb{E}}
\newcommand{\eps}{\varepsilon}
\newcommand{\dd}{\,\mathrm{d}}
\begin{document}
\section*{UnsolvedMath ID: 20000865}
\textbf{Source problem code:} catalogue entry
\subsection{Definitions and mathematical statement}
This source is a context-only Q\&A after R\"udiger von der Heydt's talk, ``Neural Mechanisms of Perceptual Inference,'' not an open mathematical problem. It reports that $50\%$ of recorded cells were in V1 and describes the selection as unbiased; response is total spikes in an $800$ ms window, and tested visual extent reached $8$ degrees. Receptive fields are mapped by moving an optimal bar until response ceases, with an explicit comment that other methods give different sizes. The speaker treats figure--ground and ownership as the same terminology, reports half-strength ownership signal at $70$ ms and IT silence up to $100$ ms, leaves an illusory-contour locality question unanswered, calls summation-field definitions controversial and summation insufficient as an explanation, and only suspects that some results persist in natural scenes.

The attached aid gives a conditional timing audit. Let $h_B$ be target half-rise time and $s_S$ the onset of a putative source signal, measured against a common baseline, and let $d_{\min}\ge0$ be a certified lower bound on every causal route lag. A sole-source account with no target response before source onset plus lag requires $h_B\ge s_S+d_{\min}$. If simultaneous uncertainty bounds are
\[
h_B\le\widehat h_B+\epsilon_B,\qquad s_S\ge\widehat s_S-\epsilon_S,
\]
then that account is excluded only by the strict certificate
\[
\widehat h_B+\epsilon_B<\widehat s_S-\epsilon_S+d_{\min}.
\]
Applied to the reported numbers, this requires $70\text{ ms}+\epsilon_B<100\text{ ms}-\epsilon_S+d_{\min}$ and a genuinely certified source-onset lower bound. Equality is insufficient. A source half-rise time is not its onset, and absence of above-threshold detected IT activity need not certify absence of a weaker causally effective source signal. The transcript does not supply these uncertainty or route-lag bounds, so its two numbers alone do not establish the certificate.

The aid also considers an inhomogeneous Poisson spike model on $[0,0.8\text{ s}]$. Its total count has law
\[
N\sim\operatorname{Poisson}\left(\int_0^{0.8\mathrm{s}}\lambda(t)\,\dd t\right).
\]
Distinct nonnegative intensity profiles with the same integral therefore give exactly the same count distribution while differing in onset or transient-versus-sustained structure. The total count alone cannot identify those dynamics in this model.

Finally, a moved-bar response profile $R(r)$ measured with threshold $\vartheta$ gives an operational radius $r_\vartheta=\sup\{r:R(r)>\vartheta\}$. This quantity depends on threshold and stimulus protocol; increasing the threshold shrinks the detected region for a monotone response profile. The aid's three-part audit distinguishes a conditional causal timing exclusion, count-based nonidentifiability and threshold-dependent receptive-field measurement. It does not turn the transcript into a solved neural-inference conjecture.
\subsection{Short English statement}
This workshop discussion gives compressed timing and measurement claims. The supplied aid specifies when timing can exclude a sole source and shows why an $800$ ms count or a threshold-defined receptive field does not uniquely identify the underlying dynamics or spatial mechanism.
\subsection{Sources}
\begin{itemize}
\item Local UnsolvedMath registry, permanent ID $20000865$, source code catalogue entry; canonical \texttt{aim-biology-notes.json}, \texttt{notes[29]}, tagged section.
\item AIM notes, \emph{Inference and Prediction in Neocortical Circuits}, PDF version dated 24 October 2003, Chapter H, ``Neural Mechanisms of Perceptual Inference,'' attributed to R\"udiger von der Heydt: \url{https://aimath.org/WWN/brain/brain.pdf}. The supplied verification expressly identifies this as Q\&A rather than an open problem. No PDF inspection is claimed here.
\item Attached research attempt 1: context only, recorded run status not a problem, rigor label proved, novelty confidence low; registry status open, classification OPEN-TRIAGE, suggested difficulty L2. It is a machine-generated, non-peer-reviewed aid requiring independent verification before citation. No independent experimental or theorem bibliography is supplied.
\end{itemize}
\end{document}
